\section{Introduction}
\label{sec:intro}

Training data rarely fits in memory in one piece, and a training loop that waits for its input wastes the hardware that it runs on. Chapter~13 of Géron's \emph{Hands-On Machine Learning}~\cite{geron} addresses this with four tools: the \texttt{tf.data} Data API, which streams, shuffles, transforms and batches data from many files at once; the TFRecord format, which stores examples as serialised Protocol Buffers; the Features API, which turns raw columns into numbers inside the model; and the TensorFlow Datasets (TFDS) project, which downloads and prepares standard benchmark data sets.

This project builds that chapter as one complete, runnable system and measures what each tool does. It has five parts, one for each question below.

\textbf{Part A.} How is a CSV data set that is split into many files read as a stream of shuffled, standardised batches, and how much does each step of that pipeline (interleaving several files, a shuffle buffer, parallel parsing, prefetching) contribute to the speed and to the quality of the shuffling?

\textbf{Part B.} How are the same rows stored as a GZIP-compressed TFRecord, how large is the file compared with the CSV, can its size be predicted exactly from the Protocol Buffers format, and does the format detect a damaged file?

\textbf{Part C.} How are a one-hot encoding of a five-category column, a 20 by 20 grid of latitude and longitude crossed into 400 cells, and learned embeddings built as part of the model, and does each of them make the predictions better?

\textbf{Part D.} How does \texttt{tfds.load} deliver a benchmark data set (MNIST) as ready batches, and how does that compare with images stored by hand in TFRecord files?

\textbf{Part E.} How is the resulting model packaged, registered, checked against quality gates and served, so that the preprocessing that was used in training is the preprocessing that is used for a request?

The data are the California housing data of Chapter~2 of the same book. A data set of this size (20,640 rows) fits in memory many times over, so the project does not claim that streaming is necessary for it. What it does is to build the pipeline that scales, to measure its behaviour on a machine with two CPU cores, and to repeat the training set up to 100 times to see how memory and speed grow when the data stop being small.

Everything is generated by one command from one configuration file. The numbers in the text, the tables and the figures of this report are written by that run: the tables are converted to \LaTeX{} by a script and the figures are drawn from the same tables, so the numbers of the tables and of the text cannot drift from the results. The numbers typed by hand are facts of the data, settings of the configuration and the run time of about 45 minutes. Section~\ref{sec:system} describes the system, Section~\ref{sec:data} the data, Sections~\ref{sec:csv} to~\ref{sec:model} the five parts, Section~\ref{sec:math} the mathematics that the experiments test, and Section~\ref{sec:limits} what is not claimed.
